\section{Описание}

В работе реализован словарь на основе собственного красно-чёрного дерева. Красно-чёрное дерево представляет собой бинарное дерево поиска, в котором каждая вершина имеет цвет, а структура поддерживает набор инвариантов, обеспечивающих логарифмическую высоту дерева и, как следствие, логарифмическое время основных операций \cite{Kormen3}.

Каждая вершина дерева хранит:
\begin{itemize}
\item строковый ключ;
\item связанное с ним значение типа \texttt{uint64\_t};
\item цвет вершины;
\item указатели на родителя, левого и правого потомка.
\end{itemize}

Для упрощения реализации пустых детей и пустого дерева используется специальная чёрная служебная вершина \texttt{nil\_}. Это позволяет не работать с \texttt{nullptr} в коде балансировки и записывать алгоритмы вставки и удаления в форме, близкой к классическому описанию \cite{Kormen3, Knuth3}.

Перед каждой операцией ключ переводится в нижний регистр, поэтому строки \texttt{word}, \texttt{Word} и \texttt{WORD} считаются одним и тем же ключом.

\subsection{Поиск}

Поиск выполняется как в обычном бинарном дереве поиска: начиная с корня, ключ сравнивается с текущей вершиной, после чего выполняется переход в левое или правое поддерево. Если достигнут \texttt{nil\_}, слово в словаре отсутствует. Время работы поиска равно $O(\log n)$.

\subsection{Вставка}

Вставка разбита на две части. Сначала новый элемент вставляется в дерево как в обычное бинарное дерево поиска. Если такой ключ уже существует, операция завершается без изменения словаря. Новый узел создаётся красным, после чего вызывается процедура \texttt{InsertFixup}, восстанавливающая инварианты красно-чёрного дерева.

В процедуре балансировки рассматриваются стандартные случаи:
\begin{itemize}
\item красный дядя -- выполняется перекраска родителя, дяди и деда;
\item чёрный дядя и \enquote{ломаная} конфигурация -- сначала выполняется поворот вокруг родителя;
\item чёрный дядя и \enquote{прямая} конфигурация -- выполняется поворот вокруг деда и перекраска.
\end{itemize}

В конце корень принудительно окрашивается в чёрный цвет.

\subsection{Удаление}

Удаление также состоит из двух этапов. Сначала вершина удаляется по правилам бинарного дерева поиска:
\begin{itemize}
\item если у неё не более одного непустого потомка, выполняется простая замена поддерева;
\item если есть оба потомка, на место удаляемой вершины переносится её inorder-последователь -- минимальная вершина в правом поддереве.
\end{itemize}

Если фактически удалённая вершина была чёрной, вызывается \texttt{DeleteFixup}. В этой процедуре рассматриваются случаи по цвету брата и его потомков. Используются перекраски и повороты, пока не будет восстановлена одинаковая чёрная высота на всех путях от текущей вершины к листьям.

\subsection{Сохранение и загрузка}

Для сохранения словаря реализован компактный бинарный формат. В файл записываются:
\begin{enumerate}
\item сигнатура формата \texttt{RBTD};
\item версия формата;
\item число элементов в словаре;
\item для каждого элемента: длина ключа, байты ключа и соответствующее значение.
\end{enumerate}

Числовые поля записываются явно в little-endian виде, поэтому формат не зависит от внутреннего представления чисел в памяти.

Загрузка выполняется во временный объект словаря. Это сделано специально: если файл повреждён, имеет неверный формат или при чтении произошла системная ошибка, текущий словарь не изменяется. Только после полного успешного чтения вызывается \texttt{Swap}, и временный словарь подменяет рабочий.

При загрузке выполняются дополнительные проверки:
\begin{itemize}
\item совпадение сигнатуры и версии;
\item корректная длина ключа;
\item допустимость символов ключа;
\item отсутствие лишних байтов в конце файла;
\item отсутствие повторяющихся ключей.
\end{itemize}

Отсутствующий или пустой файл интерпретируется как корректный пустой словарь, что соответствует условию задачи.

\subsection{Структура программы}

Программа целиком размещена в одном файле \texttt{main.cpp}. Центральной частью является класс \texttt{RedBlackTreeDictionary}, который инкапсулирует:
\begin{itemize}
\item внутреннее представление узлов;
\item операции поиска, вставки и удаления;
\item повороты и процедуры балансировки;
\item функции двоичной сериализации и десериализации.
\end{itemize}

Функция \texttt{main} отвечает только за консольный интерфейс: считывает команды, нормализует ключи, вызывает соответствующие методы словаря и печатает результат в формате, заданном условием.

Асимптотика реализованного словаря:
\begin{itemize}
\item поиск, вставка и удаление -- $O(\log n)$;
\item сохранение и загрузка -- $O(n)$;
\item дополнительная память структуры -- $O(n)$.
\end{itemize}

\section{Исходный код}

Ниже приведены ключевые фрагменты реализации, отражающие устройство словаря,
модификацию дерева и двоичную работу с файлом.

\lstinputlisting[
language=C++,
firstline=11,
lastline=97,
caption={Структура словаря и служебные объекты},
label={lst:dict-structure}
]{../main.cpp}

\lstinputlisting[
language=C++,
firstline=99,
lastline=123,
caption={Вставка элемента в дерево},
label={lst:insert}
]{../main.cpp}

\lstinputlisting[
language=C++,
firstline=125,
lastline=165,
caption={Удаление элемента из дерева},
label={lst:erase}
]{../main.cpp}

\lstinputlisting[
language=C++,
firstline=176,
lastline=225,
caption={Публичные методы сохранения и загрузки},
label={lst:save-load}
]{../main.cpp}

\lstinputlisting[
language=C++,
firstline=252,
lastline=425,
caption={Двоичная сериализация и проверка формата файла},
label={lst:file-format}
]{../main.cpp}
\pagebreak

\section{Консоль}
\begin{alltt}
$ g++ -std=c++17 -O2 main.cpp -o lab2.out
$ ./lab2.out
+ a 1
OK
+ A 2
Exist
+ zebra 17
OK
zEbRa
OK: 17
! Save dict.bin
OK
-\ zebra
OK
! Load dict.bin
OK
zebra
OK: 17
\end{alltt}
\pagebreak
